\uuid{dOzb}
\titre{Somme de variables indépendantes ou multiple d'une variable}
\auteur{Maxime Nguyen}
\organisation{AMSCC}
\datecreate{2026-11-27}
\niveau{L2}
\module{Probabilité}
\chapitre{Probabilité continue}
\sousChapitre{Loi normale}
\difficulte{2}
\theme{loi normale, somme de variables indépendantes, transformation affine}

\contenu{
\begin{question}
On note $\mathcal{N}(\mu,\sigma)$ la loi normale d'espérance $\mu$ et d'écart-type $\sigma$. Soit $X_1$, $X_2$, $X_3$ trois variables indépendantes de même loi $\mathcal{N}(\mu, \sigma)$. Quelles affirmations sont vraies ?
\begin{answers}
\bad{$X_1 + X_2 + X_3$ suit la loi $\mathcal{N}(3\mu, 3\sigma)$.}
\bad{$X_1 + X_2 + X_3$ et $3X_1$ ont la même loi.}
\good{$3X_1$ suit la loi $\mathcal{N}(3\mu, 3\sigma)$.}
\bad{$\mathbb{V}(3X_1) = 3\,\mathbb{V}(X_1)$.}
\good{$X_1 + X_2 + X_3$ suit la loi $\mathcal{N}(3\mu, \sqrt3\,\sigma)$.}
\end{answers}
\begin{explanations}
E : par indépendance, $\mathbb{V}(X_1 + X_2 + X_3) = 3\sigma^2$, d'où l'écart-type $\sqrt3\,\sigma$.

C : $\mathbb{V}(3X_1) = 9\sigma^2$, d'où l'écart-type $3\sigma$ (et D est faux).

A, B sont donc faux. Additionner trois mesures indépendantes ne triple pas la dispersion : les écarts se compensent en partie. Tripler une seule mesure triple son écart.

Contre-sens visé : confondre $X_1 + X_2 + X_3$ (variables indépendantes)
et $3X_1$ (une seule variable recopiée trois fois).
\end{explanations}
\end{question}
}
